\section{hitchtrace}
\label{sec:design}

\sys explains over-budget frames while emitting nothing for normal frames,
with no change to the game or the kernel. \autoref{fig:record} is its
record for the frame of \autoref{sec:example}: the 30 waits for the server
form one group of 8.9\,ms, and its chain says that the server was kept off
the CPU for 5.5\,ms, queued behind a thread of the background program.

\begin{figure*}[t]
\centering
\begin{lrbox}{\recordbox}
\footnotesize
\begin{minipage}{1.42\textwidth}
\begin{verbatim}
hitch frame 2475  23.5 ms  (budget 20.8 ms, over by 2.7 ms, 852% kernel-observable)  root mg_main/39791  [open stall]
  oncpu      0.6 ms   runqueue 1.7 ms
  blocked   21.1 ms  futex 8.1 | timer 4.0 | present 0.2 | task 8.9
  resolved   waiting on a running task 8.9 ms, inherited stalls 8.2 ms
  stalls     off-CPU time by cause and by who ended it, largest first
      8.9 ms  task     blocked in __skb_wait_for_more_packets, 30 times, longest 5.5 ms
              <- mg_server/39784        kept off the CPU 5.5 ms, last queued behind bgload2/39824
      8.1 ms  futex    blocked in __futex_wait
              <- mg_worker1/39806       was running (waiting for work) 0.1 ms
      4.0 ms  timer    blocked in do_nanosleep
              <- [hardirq]              woken by an interrupt
      1.6 ms  runqueue waiting for a CPU, 31 times, longest 0.2 ms
      0.3 ms  in 2 smaller groups
  threads
    mg_main/39791          (root)  oncpu 0.6  runqueue 1.7  blocked 21.1 futex 8.1 | timer 4.0 | present 0.2 | task 8.9
    mg_worker2/39807               oncpu 6.2  runqueue 0.9  blocked 15.2 futex (still blocked at frame end 1.2 ms)   largest 16.1 ms in __futex_wait
    mg_worker1/39806               oncpu 8.1  runqueue 0.1  blocked 15.2 futex (still blocked at frame end 0.1 ms)   largest 15.3 ms in __futex_wait
    mg_worker0/39805               oncpu 7.1  runqueue 0.0  blocked 15.2 futex (still blocked at frame end 1.1 ms)   largest 15.2 ms in __futex_wait
    [3 more threads]
  34 stalls this frame, 45 off-CPU intervals of the process logged
\end{verbatim}
\end{minipage}
\end{lrbox}

\resizebox{\textwidth}{!}{\usebox{\recordbox}}
\caption{A \sys record, as printed, for the frame of \autoref{sec:example};
three thread lines are cut. The game's main thread made 30 blocking calls
to a server process in this frame; two were slow because the server, woken,
waited for its CPU behind a thread of a background program. The header's
share (852\%) is what the prototype's gate weighs: it counts the frame
limiter's sleep and the wait for workers that were running, which is why
the resolved gate subtracts the latter (\autoref{sec:eval-gate}). The same
record, with its interval log, is written as JSONL.}
\label{fig:record}
\end{figure*}

\subsection{Design Overview}

\autoref{fig:arch} shows the components. Hooks on scheduler, syscall,
fault and interrupt events keep per-task state in task-local storage
(\autoref{sec:ebpf}); a uprobe at each frame mark decides whether to emit a
record through the ring buffer.

\begin{figure*}[t]
\centering
\resizebox{\textwidth}{!}{\begin{tikzpicture}[
  font=\footnotesize,
  box/.style={draw, rounded corners=2pt, align=center, inner sep=3pt, minimum height=0.62cm},
  hook/.style={box, fill=blue!6, minimum width=2.35cm},
  state/.style={box, fill=orange!10},
  prog/.style={box, fill=green!8},
  zone/.style={draw, dashed, rounded corners=4pt, inner sep=6pt, black!55},
  arr/.style={-{Stealth[length=1.6mm]}, thick},
  zlabel/.style={font=\scriptsize\itshape, text=black!60, anchor=south west, inner sep=1pt}
]
% ---- game process
\node[box] (threads) at (1.1,0) {game threads\\(render, workers, \ldots)};
\node[box, right=0.45cm of threads] (loader) {Vulkan\\loader};
\node[box, fill=yellow!15, right=0.45cm of loader, font=\scriptsize] (layer) {\textbf{\footnotesize implicit layer} \texttt{VK\_LAYER\_hitchtrace}\\\texttt{hitch\_frame\_mark()} $\to$ next \texttt{vkQueuePresentKHR} $\to$ \texttt{hitch\_present\_end()}};
\node[box, right=0.45cm of layer] (icd) {driver\\(ICD)};
\draw[arr] (threads) -- (loader);
\draw[arr] (loader) -- (layer);
\draw[arr] (layer) -- (icd);
\node[zone, fit=(threads)(icd)(layer)] (game) {};
\node[zlabel] at (game.north west) {game process (unmodified; native or Proton)};

% ---- hooks
\node[hook] (h1) at (1.3,-1.9) {\texttt{sched\_switch}};
\node[hook] (h2) at (3.8,-1.9) {\texttt{sched\_waking}};
\node[hook] (h3) at (6.3,-1.9) {\texttt{sys\_enter/exit}};
\node[hook] (h4) at (8.8,-1.9) {faults, reclaim,\\compaction};
\node[hook] (h5) at (11.3,-1.9) {GPU fence /\\syncobj waits};
\node[hook] (h6) at (13.8,-1.9) {\texttt{irq\_exit\_rcu}\\entry/exit};

% ---- task storage spans all hooks
\node[state, minimum width=14.9cm, text width=14.5cm] (ts) at (7.55,-3.35)
  {\textbf{task-local storage, every task}: switch-out/in times $\cdot$ state at switch-out $\cdot$ syscall class $\cdot$ time in faults, reclaim, compaction $\cdot$ \textbf{off-CPU history} (32 checkpoints) $\cdot$ last CPU holder $\cdot$ \textbf{wake links}: the waker's off-CPU time during the wait and who held its CPU, and the waker's own links, up to 5 hops};
\foreach \h in {h1,h2,h3,h4,h5,h6} \draw[arr] (\h.south) -- (\h.south |- ts.north);

% ---- per-frame state
\node[state, text width=4.6cm] (root) at (5.2,-4.85) {\textbf{root timeline}: 15 cause buckets partitioning the frame; resolved inherited / wait-on-running; \textbf{stall groups}, each with its chain};
\node[state, text width=3.0cm] (slots) at (9.6,-4.85) {\textbf{per-thread slots}: each thread's buckets, longest stall, preemptor};
\node[state, text width=2.6cm] (evlog) at (13.2,-4.85) {\textbf{interval log}: every off-CPU interval of the process, two buffers};
\draw[arr] (root.north |- ts.south) -- (root.north);
\draw[arr] (slots.north |- ts.south) -- (slots.north);
\draw[arr] (evlog.north |- ts.south) -- (evlog.north);

% ---- frame mark program
\node[prog, text width=6.0cm] (mark) at (7.3,-6.35) {\textbf{frame mark} (uprobe on present entry): close frame $\to$ budget and cause gate $\to$ fill record, or discard $\to$ open next frame};
\draw[arr] (root.south) -- (root.south |- mark.north);
\draw[arr] (slots.south) -- (slots.south |- mark.north);
\draw[arr] (evlog.south) -- (evlog.south |- 0,-6.15) -- ([yshift=0.2cm]mark.east);
\draw[arr, dashed] (layer.south) -- (layer.south |- 0,-1.05) -- (15.55,-1.05) -- (15.55,-6.55) -- ([yshift=-0.2cm]mark.east);
\node[font=\scriptsize\itshape, text=black!60, anchor=south east] at (15.5,-1.05) {uprobes on both markers};

\node[box, fill=gray!10] (rb) at (1.6,-6.35) {ring buffer\\(over-budget frames only)};
\draw[arr] (mark.west) -- (rb.east);
\node[zone, fit=(h1)(h6)(ts)(root)(slots)(mark)(rb)] (kern) {};
\node[zlabel] at (kern.north west) {kernel (eBPF, verified)};

% ---- user space
\node[box] (fe) at (2.2,-8.1) {loader and ring-buffer consumer};
\node[box] (sym) at (6.6,-8.1) {kernel/user stack symbolizer};
\node[box] (out) at (10.6,-8.1) {human report and JSONL};
\draw[arr] (rb.south) -- (rb.south |- fe.north);
\draw[arr] (fe) -- (sym);
\draw[arr] (sym) -- (out);
\node[zone, fit=(fe)(out)] (usr) {};
\node[zlabel] at (usr.north west) {\sys (user space)};
\end{tikzpicture}
}
\caption{Overview of \sys. Hooks keep per-task state current as the
scheduler runs; the frame mark closes the root thread's frame, evaluates the
gate, and emits a record, followed by the frame's interval log, only for an
over-budget frame.}
\label{fig:arch}
\end{figure*}

Four principles answer the challenges of \autoref{sec:challenges}.

\paragraph{Account on the root, not the process (C1, C3)}
Summing off-CPU time over all of a game's threads is unbounded: a process
with ten parked workers accumulates ten frames' worth of off-CPU time per
frame. \sys accounts on the root, whose states partition the frame time;
other threads enter the record as context and through the chain.

\paragraph{Group waits, then explain each group (C1)}
The root's waits are summed by cause and by who ended them, and each group
is explained through the chain of its longest wait, so forty short waits on
one thing count as one stall.

\paragraph{Build chains forward, not backward (C2, C3)}
Instead of recording wakeup edges and searching them when a frame goes bad,
every task carries the chain that led to its latest wakeup, copied forward
from its waker at wakeup time, with how long that waker was itself off-CPU.
Resolving a wait is a lookup, not a search, and it tells a stalled waker
from a busy one.

\paragraph{Decide at frame end, in the kernel (C1, C4)}
All per-frame state lives in maps until the frame mark. Only then does \sys
know whether the frame was over budget, and only then does it reserve
ring-buffer space.

\subsection{Frame Source}
\label{sec:layer}

\sys needs two points per frame: present entry, which ends one frame and
begins the next, and present return, which ends the present call. The
implicit layer \texttt{VK\_LAYER\_hitchtrace} wraps
\texttt{vkQueuePresentKHR} in the device dispatch chain and calls two
exported, non-inlined marker functions around the next layer's present,
passing a process-global present counter as the frame id. \sys attaches a
uprobe to each marker in the layer's shared object. The layer is inert
unless \texttt{HITCHTRACE=1} is set; with no probe attached the markers are
two empty calls. A uprobe on the Vulkan loader's exported
\texttt{vkQueuePresentKHR} would not do (C5): applications that resolve
device functions through \texttt{vkGetDeviceProcAddr}, as volk, DXVK,
vkd3d-proton and Wine's Vulkan bridge do, receive a pointer into the first
layer or the driver and never call the export~\cite{vulkan-loader}; on our
machine \texttt{vkcube} imports no \texttt{vk*} symbol at all. The layer
sits in the dispatch chain however the application obtained its pointers,
and because Steam's container runtime resolves a host layer to the same
inode, a uprobe attached through the host path fires inside it.
\todo{Confirm inside pressure-vessel on the evaluation machine; the layer
is 64-bit only.}

Programs without a layer can name any symbol as the frame marker.

\paragraph{Frame model} Frame $N$ runs from present entry $N-1$ to present entry
$N$. Blocked time of the root between present entry and present return is
charged to \bucket{block\_present}, the display pacing the application, so a
vsynced game does not register its vsync waits as stalls. The bracket is
only opened when the return probe is attached; otherwise a missed close
would silently relabel every wait.

\subsection{Root Timeline and Cause Taxonomy}
\label{sec:taxonomy}

Within a frame $[s, e)$, the root is always in one of four states, and every
interval of each state is clipped to the window:
\begin{equation*}
\textit{oncpu} + \textit{runnable} + \textit{runqueue} + \textstyle\sum \textit{blocked}_c = e - s .
\end{equation*}
On-CPU time is measured from switch-in to switch-out, never computed as a
remainder. At switch-out, \sys records the time, whether the switch was a
preemption, the task's \texttt{in\_iowait} bit, and a kernel stack for
voluntary sleeps. At switch-in, the interval since switch-out is split at
the wakeup timestamp: before the wakeup is blocked time, after it is
run-queue delay. \autoref{tab:taxonomy} lists the buckets and how each is
decided.

\begin{table}[t]
\centering
\footnotesize
\newcolumntype{L}[1]{>{\raggedright\arraybackslash}p{#1}}
\setlength{\tabcolsep}{3pt}
\begin{tabular}{@{}L{1.85cm}L{2.75cm}L{2.95cm}@{}}
\toprule
\textbf{Bucket} & \textbf{Hook} & \textbf{Rule} \\
\midrule
\bucket{oncpu} & \texttt{sched\_switch} & running, minus the carve-outs below \\
\bucket{oncpu\_fault}, \bucket{\ldots\_major} & fentry/fexit \texttt{handle\_\allowbreak mm\_\allowbreak fault} & on-CPU inside a fault; major from \texttt{VM\_FAULT\_\allowbreak MAJOR} \\
\bucket{oncpu\_\allowbreak reclaim} & direct and memcg reclaim tracepoints & on-CPU inside reclaim \\
\bucket{oncpu\_\allowbreak compact} & compaction tracepoints & on-CPU inside compaction \\
\midrule
\bucket{runnable} & \texttt{sched\_switch} & switched out while runnable; preemptor recorded \\
\bucket{runqueue} & \texttt{sched\_waking} & woken, waiting for a CPU \\
\midrule
\bucket{block\_\allowbreak present} & layer markers & blocked inside present \\
\bucket{block\_gpu} & \texttt{dma\_fence\_\allowbreak wait\_\allowbreak timeout}, \texttt{drm\_syncobj\_\allowbreak *wait\_\allowbreak ioctl} & blocked in a GPU fence wait \\
\bucket{block\_io} & \texttt{in\_iowait} at switch-out & blocked in I/O wait \\
\bucket{block\_futex} & \texttt{sys\_enter/\allowbreak exit} & in \texttt{futex}, \texttt{futex\_\allowbreak waitv}, \texttt{futex\_\allowbreak wait} \\
\bucket{block\_poll} & \texttt{sys\_enter/\allowbreak exit} & in poll, select, epoll \\
\bucket{block\_timer} & \texttt{sys\_enter/\allowbreak exit}, waker & in a sleep, or woken by an interrupt \\
\bucket{block\_task} & \texttt{sched\_waking} & woken by another task, none of the above \\
\bucket{block\_other} & --- & none of the above \\
\bottomrule
\end{tabular}

\caption{Root-timeline buckets, in the precedence order used for blocked
time. Together they partition the frame.}
\label{tab:taxonomy}
\end{table}

\paragraph{What was asked beats who answered} A blocked interval is named first
by what the thread asked the kernel for, and only then by who woke it: a
futex wait ended by its timeout interrupt is still a futex wait. The
present bracket outranks everything, including GPU fences, because a FIFO
present may legitimately block on a fence and blaming the GPU would point at
the application's own rendering.

\paragraph{On-CPU kernel stalls} Page faults, reclaim and compaction are mostly
on-CPU. \sys times them per task with entry/exit hooks and, at each
switch-out and at frame end, carves the time accrued since the last update
out of \bucket{oncpu}. Reclaim is hooked at both the direct and the memcg
tracepoints, since a cgroup memory limit triggers only the latter. Reclaim
and compaction usually run inside the fault that needs the memory, so the
fault and the reclaim around a compaction subtract what ran nested inside
them, and the three carve up the same on-CPU time once.

\paragraph{Per-thread context} Every thread of the target process also keeps a
per-frame line in one of 32 slots: its own buckets, its longest off-CPU
interval, the task that preempted it longest, and, if it is still off-CPU
when the frame closes, how long it has been waiting. These lines are context:
threads run in parallel, so their off-CPU time does not add to the frame.

\subsection{Online Wake-Chain Resolution}
\label{sec:chain}

\begin{figure}[t]
\centering
\begin{tikzpicture}[
  font=\scriptsize, x=0.7cm, y=0.8cm,
  arr/.style={-{Stealth[length=1.5mm]}, thick}
]
\def\lane#1#2{\node[anchor=east, font=\scriptsize\bfseries] at (-0.1,#1) {#2};}
\lane{2}{render}
\lane{1}{worker}
\lane{0.1}{IRQ}
% render (root)
\fill[cOncpu!70] (0,1.8) rectangle (1.5,2.2);
\fill[cFutex!55] (1.5,1.8) rectangle (7,2.2);
\fill[cSched!45] (7,1.8) rectangle (7.5,2.2);
\fill[cOncpu!70] (7.5,1.8) rectangle (10,2.2);
\node at (4.25,2) {blocked: futex};
% worker
\fill[cOncpu!70] (0,0.8) rectangle (2.5,1.2);
\fill[cIO!45] (2.5,0.8) rectangle (5.4,1.2);
\fill[cSched!45] (5.4,0.8) rectangle (5.8,1.2);
\fill[cOncpu!70] (5.8,0.8) rectangle (7,1.2);
\fill[cFutex!25] (7,0.8) rectangle (10,1.2);
\node at (3.95,1) {blocked: io};
% wakeups
\draw[arr] (5.4,0.15) -- (5.4,0.78);
\node[anchor=west] at (5.5,0.2) {disk IRQ wakes worker};
\draw[arr] (7,1.22) -- (7,1.78);
\node[anchor=east] at (6.9,1.5) {worker wakes render};
% frame window
\draw[black!55, dashed] (0.3,-0.25) -- (0.3,2.3);
\draw[black!55, dashed] (9.7,-0.25) -- (9.7,2.3);
\node[font=\tiny, text=black!65, anchor=north] at (0.3,-0.3) {frame start};
\node[font=\tiny, text=black!65, anchor=north] at (9.7,-0.3) {frame mark};
% resolution of render's blocked interval [1.5, 7)
\draw[decorate, decoration={brace, amplitude=3pt}] (1.52,2.3) -- (2.48,2.3)
  node[midway, above=3pt, align=center, font=\tiny] {wait on\\running};
\draw[decorate, decoration={brace, amplitude=3pt}] (2.52,2.3) -- (5.78,2.3)
  node[midway, above=3pt, align=center, font=\tiny] {inherited:\\worker's io + run queue};
\draw[decorate, decoration={brace, amplitude=3pt}] (5.82,2.3) -- (6.98,2.3)
  node[midway, above=3pt, align=center, font=\tiny] {wait on\\running};
\end{tikzpicture}

\caption{Resolving a blocked interval of the root. The part of render's
futex wait during which the worker was itself off-CPU (I/O plus run queue,
ended by the disk interrupt) is inherited; the rest is time the worker spent
running, i.e.\ render waiting for work. The reported chain is
render $\leftarrow$ worker (io) $\leftarrow$ [irq].}
\label{fig:chain}
\end{figure}

\paragraph{Copy-forward chains} On every \texttt{sched\_waking}, \sys writes into
the wakee's task storage a \emph{link} describing the waker as it is at that
moment: its identity, how long it has itself been off-CPU since the wakee
began to wait, split into runnable and blocked time, the cause and kernel
stack of its last blocked interval, and the task that last held its CPU in
that window. It then copies the waker's own links into the wakee's
slots $1..4$, so the wakee holds its waker, its waker's waker, and so on, up
to five hops. The copy has constant size, as the verifier requires, and a
sequence counter on the waker detects a copy torn by a concurrent wakeup of
the waker on another CPU; a torn copy is discarded rather than trusted.

\paragraph{Off-CPU history} A waker's time off-CPU during someone else's wait is
rarely one interval: a worker that shares its CPU with another process loses
it several times a frame, and a read split into small requests blocks many
times. Every task therefore keeps a ring of 32 checkpoints of its
cumulative runnable and blocked time, each taken just before an off-CPU
interval began. Intervals that end within 2\,ms of the newest checkpoint
share it rather than open their own, so the ring reaches back at least
64\,ms however finely the time is chopped up, at the cost of up to one grain
at the start of a window that opens among merged intervals. At a wakeup the
waker reads its off-CPU time since the wakee began to wait from the
checkpoint before that moment. A history shorter than the wait is flagged,
and the hop is then a lower bound.

\paragraph{Resolution} When the root switches back in after a blocked interval
$[\textit{lo}, \textit{hi})$ ended by another task,
\begin{align*}
\textit{inherited} &= \min\bigl(\textit{hi} - \textit{lo},\; r_w + b_w\bigr),\\
\textit{wait\_oncpu} &= (\textit{hi} - \textit{lo}) - \textit{inherited},
\end{align*}
where $r_w$ and $b_w$ are the waker's runnable and blocked time during the
wait, from its link (\autoref{fig:chain}). Inherited time is a stall
one hop away; \textit{wait\_oncpu} is the root waiting for a thread that was
running, the game's own work rather than a kernel stall. Both are reported
alongside the partition. For the longest interval of each stall group
(\autoref{sec:groups}), \sys repeats the step hop by hop down the chain, each
hop's window narrowed to the previous hop's wait, recording for each hop the
task, the time it explains and how much of it was runnable, what it was
blocked in, and whether its wait had ended before the previous hop's began
(a stale link).

\paragraph{Where chains stop} A chain ends at an interrupt, at the idle task, at
a waker that was running (off-CPU for less than the minimum stall,
100\,$\mu$s), or at a waker that spent most of the wait runnable rather than
blocked. In the last case the cause is whoever held its CPU, not whoever had
last woken it, and the hop names that holder instead of continuing.
\texttt{sched\_waking} runs in the waker's context, but a wakeup from an
interrupt runs on whichever task the interrupt landed on; taking that task
as the waker would splice a bystander's history into the chain. \sys
detects hardirq and NMI context from the
per-CPU preemption count, read through two weak kernel symbols because x86
moved the variable into and back out of a per-CPU structure between kernel
versions; the verifier prunes the branch for the symbol that does not exist.
Softirqs that run on interrupt exit, after the hardirq bits are cleared, are
detected by a per-CPU nesting counter maintained by \texttt{fentry} and
\texttt{fexit} on \texttt{irq\_exit\_rcu} and reset at every context switch,
where it must be zero. An interrupt link carries no task identity: the task
it names is a bystander.

\paragraph{Who held the CPU} A task loses its CPU in one of two ways. Preempted,
it is switched out while runnable, and \texttt{sched\_switch} names the task
that took it. Woken while its CPU is busy, it waits in the run queue without
ever being preempted, and no event names anyone; at the switch-in that ends
a run-queue wait of at least 100\,$\mu$s, \sys records the task leaving the
CPU as the last holder, flagged \emph{queued}. Both are kept for every task,
since any task can be waited for, and the link carries the waker's last
holder in the wait window. Only the last holder is seen: not an earlier one
in the same wait, and not one on a CPU the task migrated from.

\subsection{Stall Groups}
\label{sec:groups}

Explaining only the frame's longest interval picks the wrong wait whenever
the time is lost in many short ones (C1), and can never pick preemption,
which is not a blocked interval. \sys therefore groups the root's off-CPU
intervals within the frame by cause and by who ended them: blocked intervals
by the thread that woke the root (or by an interrupt), time spent preempted
by the process that took the CPU, and time spent woken but queued for a CPU
in a group of its own. Each group keeps its total, its count, its longest
interval, its inherited time, and the chain behind its longest interval.
Seven groups get a slot. An interval of at least 100\,$\mu$s takes a free
slot, or displaces the smallest group once it is longer than that group's
whole total; whatever is displaced or too short to be a stall is summed in
an eighth group, so the groups always add up to the root's off-CPU time.
The front end ranks the groups by total and prints each one above 3\% of
the frame with its chain (\autoref{fig:record}).

\subsection{Cause Gate and Emission}
\label{sec:gate}

At the frame mark, \sys closes the root's on-CPU tail, applies the
carve-outs, and computes the frame time $f$ and the budget $b$. The budget is
fixed, or, for an uncapped game, 1.5 times the rolling median of the last 19
frames, with a floor. The frame is a hitch if $f > b$. With a share $\sigma
> 0$, the frame is emitted only if
\[
\textit{stall} \geq \sigma \cdot (f - b),
\]
where \textit{stall} is the root's time in every bucket except
\bucket{oncpu} and \bucket{block\_present}. A share of zero gates on the
symptom alone, which is what an unbiased prevalence study needs; a share of
0.5 says ``at least half of the overrun must be stalls.''

For an emitted frame, \sys reserves a 10,344-byte record in the ring buffer
and fills it with the partition, the resolution totals, the stall groups
with their chains, and the 32 per-thread lines. The record is followed by
the frame's \emph{interval log}: every off-CPU interval of every thread of
the process during the frame, saying when the thread left the CPU, when it
was woken and when it ran again, why it waited, who woke or preempted it,
how long that waker had itself been off-CPU, and where it blocked. Each
interval is appended to a per-frame buffer as it ends, at 72 bytes, a median
of 16 to 22 per hitch in our case-study game (\autoref{sec:eval-cases}). For
every other frame \sys empties the buffer, resets the root's state and moves
on: no bytes cross into user space. The user-space front end symbolizes
kernel stacks, optionally user stacks, writes each kernel stack the record
refers to once in a table at its end, and produces a human-readable report
and JSONL (\autoref{fig:record}). A viewer turns the JSONL into a page
that draws each hitch's intervals as one lane per thread, with an arrow from
whichever thread ended each wait, beside its groups and chains.
